% !TEX TS-program = pdflatex
% !TEX encoding = UTF-8 Unicode

% This is a simple template for a LaTeX document using the "article" class.
% See "book", "report", "letter" for other types of document.

\documentclass[12pt]{book} % University guidelines recommend 12pt.
\usepackage[utf8]{inputenc} % set input encoding (not needed with XeLaTeX)

%%% Examples of Article customizations
% These packages are optional, depending whether you want the features they provide.
% See the LaTeX Companion or other references for full information.

%%% PAGE DIMENSIONS
\usepackage{geometry} % to change the page dimensions
\geometry{a4paper} % or letterpaper (US) or a5paper or....
\geometry{margin=3cm} % for example, change the margins to 2 inches all round
% \geometry{landscape} % set up the page for landscape
%   read geometry.pdf for detailed page layout information

\usepackage{graphicx} % support the \includegraphics command and options
\graphicspath{{Figures/}} % Set the default folder for images

% \usepackage[parfill]{parskip} % Activate to begin paragraphs with an empty line rather than an indent

%%% PACKAGES
\usepackage{indentfirst}
\usepackage{booktabs} % for much better looking tables
\usepackage{array} % for better arrays (eg matrices) in maths
\usepackage{paralist} % very flexible & customisable lists (eg. enumerate/itemize, etc.)
\usepackage{verbatim} % adds environment for commenting out b6locks of text & for better verbatim
%\usepackage{subfigure} % make it possible to include more than one captioned figure/table in a single float
\usepackage{cleveref}  % For referencing more than one equation at at time
\usepackage[colorlinks=true,urlcolor=darkblue,citecolor=blue,linkcolor=black]{hyperref} % for hyper referencing
\renewcommand{\sectionautorefname}{\S}  % To reference sections using a fancy symbol.
\usepackage{amsmath}
\usepackage[numbers,sort&compress]{natbib}  % Make a group of citations appear in order in the document
\usepackage{bbold}
\usepackage[toc,page]{appendix}
\usepackage[labelfont=bf]{subcaption}
\usepackage[labelfont=bf]{caption}
\captionsetup{subrefformat=parens}
% These packages are all incorporated in the memoir class to one degree or another...
\usepackage{wrapfig}% http://ctan.org/pkg/wrapfig
\usepackage{lipsum}% http://ctan.org/pkg/lipsum
\usepackage{epstopdf}
\epstopdfsetup{update}

%%% HEADERS & FOOTERS
\usepackage{fancyhdr} % This should be set AFTER setting up the page geometry
\pagestyle{fancy} % options: empty , plain , fancy
\renewcommand{\headrulewidth}{0pt} % customise the layout...
\lhead{}\chead{}\rhead{}
\lfoot{}\cfoot{\thepage}\rfoot{}
\usepackage{epigraph}
\setlength\epigraphwidth{.9\textwidth}
%\newcommand{\epigraphflush}{flushright}
\setlength\epigraphrule{1pt}

%%% SECTION TITLE APPEARANCE
\usepackage{sectsty}
\allsectionsfont{\sffamily\mdseries\upshape} % (See the fntguide.pdf for font help)
% (This matches ConTeXt defaults)

% New commands
\newcommand\numberthis{\addtocounter{equation}{1}\tag{\theequation}}  % For numbering only one equation in an array.

%%% ToC (table of contents) APPEARANCE
\usepackage[nottoc,notlof,notlot]{tocbibind} % Put the bibliography in the ToC
\usepackage[titles,subfigure]{tocloft} % Alter the style of the Table of Contents
\renewcommand{\cftsecfont}{\rmfamily\mdseries\upshape}
\renewcommand{\cftsecpagefont}{\rmfamily\mdseries\upshape} % No bold!
%%% END Article customizations

%%% The "real" document content comes below...
% \title{Brownian motors thermally coupled to the environment}
% \author{Jack Devine}
%\date{} % Activate to display a given date or no date (if empty),
         % otherwise the current date is printed

\begin{document}
\frontmatter
	\begin{titlepage}
		\centering
		{\huge\bfseries Self-induced temperature gradients in Brownian dynamics\par}
		\vspace{1cm}
		{\huge Jack Devine\par}
		\vspace{1cm}
		{\scshape\LARGE The University of Otago \par}
		\vspace{1cm}
		\includegraphics[width=0.24\textwidth]{OtagoLogo.eps}\par\vspace{1cm}
		\vspace{1cm}
		{\scshape\Large A thesis submitted for the partial fulfillment of  Bachelor of Science with Honours in physics at the University of Otago, Dunedin, New Zealand\par}
		\vspace{2cm}
		\vfill
		supervised by\par
		Dr M. W. \textsc{Jack}

		\vfill

	% Bottom of the page
		{\large \today\par}
	\end{titlepage}


\newpage
%----------------------------------------------------------------------------------------
%	ABSTRACT
%----------------------------------------------------------------------------------------
%
\begin{center}
\section*{Abstract} % This section will not appear in the table of contents due to the star (\section*)
As Brownian particles diffuse over a potential landscape, they draw energy from their environment in the form of heat. We account for these interactions by coupling the Smoluchowski equation to a modified version of the heat equation yielding two coupled non-linear partial differential equations. We show that the interactions can induce temperature gradients that significantly alter the dynamical behavior of the particle. We obtain theoretical steady state solutions to the equations of motion and employ numerical techniques to explore the dynamical properties of the system. We find the Kramers' rate for our modified view of Brownian dynamics and we make concrete observations of Brownian cooling and the reverse Landauer blowtorch effect. We show that these phenomena can significantly alter the motion of a particle on a tilted periodic potential, in particular we show that for particular sets of parameters, the movement of the particle can be brought to a stand-still. This has direct consequences to biology in the context of molecular motors as well as the dynamics of colloidal particles.
\end{center}

\clearpage

\section*{Acknowledgments}
\noindent First I would like to Acknowledge my supervisor Dr. Michael Jack for helping me through this difficult project and always being available to give me advice. As well as this, I would like to acknowledge Professor Michael Colombo, Vaughan Weatheral, Luke Trainor and Loren Kersey for reading very early drafts of this report, they helped me to remove many mistakes. Finally, I would like to thank my colleagues in the honors room for their camaraderie and banter throughout the year.

\clearpage

\null\vfill
{
\centering{
{\Large%
  At the atomic level, we have new kinds of forces and new kinds of possibilities, new kinds of effects. The problems of manufacture and reproduction of materials will be quite different. I am, as I said, inspired by the biological phenomena in which chemical forces are used in a repetitious fashion to produce all kinds of weird effects (one of which is the author).} \\
\vspace{1cm}
  \raggedleft{\Large\textbf{ Richard Feynman \\ -- There's Plenty of room at the bottom (1959).}}\par%
}
}
\vfill\vfill

\clearpage

%----------------------------------------------------------------------------------------

\newpage % Start the article content on the second page, remove this if you have a longer abstract that goes onto the second page

\tableofcontents



%----------------------------------------------------------------------------------------
%	INTRODUCTION
%----------------------------------------------------------------------------------------

\mainmatter
\chapter{Introduction}

Brownian motors are devices that can use stored energy to create directed motion on a microscopic scale. As well as being able to crank a rotor in the in the fashion of a traditional motor, they are also able to pump ions against a gradient and translocate molecules. Brownian motors have been implemented in the laboratory, for example Ref \cite{BlickleBechinger2011} created a stochastic heat engine by placing a single colloidal particle in a time dependent optical trap. Likewise, Ref \cite{Pedro2014} placed a colloidal particle in an optical tweezer and drove the particle with explosive vaporization of the surrounding liquid, thus demonstrating a thermal mechanism for Brownian motors. Brownian motors also include devices that can transport molecules over a long distance, for example Ref \cite{JoelBader1999} placed DNA molecules in a time dependent potential to transport the molecules. A very important class of Brownian motors are those where the energy is supplied by a chemical reaction. None of the above experiments fit this criteria, however these types of motors are ubiquitous in biology \cite{PhillipsQuakeMay2006, Magnasco1994}. Recently, thanks to improvements in imaging techniques, researchers have been able to make highly detailed images of these motors and their working components \cite{YiWeiChang2016}.

In this project, we will model Brownian motors using concepts from statistical mechanics. In particular, a Brownian motor will be modeled by a Brownian particle diffusing over its free energy landscape \cite{Reimann2001}, which in one dimension will be denoted by $x$. In the case of Brownian motion, it is natural to think of $x$ as a spatial coordinate, however in the case of Brownian motors this is not always the case. Often we will think of $x$ as a reaction coordinate for a chemical reaction, or in the case of a rotary motor, it could be the angle of the motor. We will think of a Brownian particle moving in titled periodic potential of the form $V(x) = f x + v_0(x)$ for some external forcing $f$ and some periodic function $v_0(x)$ with period $L$. This is shown schematically in figure \ref{fig:Schematic}, in this figure we have a particle with a certain known probability density. The particle is agitated by thermal vibrations in a random diffusive manner, however there is also a forcing on the particles that we describe using the potential. Different types of Brownian motors have been explored in the literature, including the Feynman ratchet \cite{Feynman1963}, the Landauer blowtorch \cite{Landauer1988}, thermal ratchets \cite{Pedro2014}, time dependent potentials \cite{JoelBader1999,BlickleBechinger2011} and tilted periodic potentials \cite{Leibler1993,Magnasco1994}.

\begin{figure}[tb]
	\centering
	\subfigure{%
		\includegraphics[width=0.45\columnwidth]{SchematicConstantTempInit}
	}
\quad
	\subfigure{
		\includegraphics[width=0.45\columnwidth]{SchematicConstantTempFinal}
	}
\caption{Schematic showing the probability density of particles diffusing in a one dimensional periodic potential at a fixed temperature. We see that the particles tend to drift down the potential as they diffuse, this drift will be called the current $J$ which we will quantify in section \ref{Smoluchowski}.}
\label{fig:Schematic}
\end{figure}

\section{Classes of Brownian motors} \label{BrownianMotorClasses}
Here we will discuss different classes of Brownian motors and their relationship to this project.

\subsection{Feynman ratchet and pawl}
The Feynman ratchet was initially discussed in the Feynman lectures \cite{Feynman1963} and was at first thought to be able to achieve greater than Carnot efficiency, however closer analysis showed that this was not possible \cite{ParrondoEspanol1996}. The system works as follows, we have two boxes that are thermally insulated from one another that are connected by an axle that can rotate. In one box there is a ratchet and pawl connected to the axle that makes it easy for the axle to turn one way (say clockwise), but hard to turn the other way (anti-clockwise). In the other box the axle is connected to paddles that are being buffeted by a gas (which we call the bath). The motion of the paddles are random since they are dictated by Brownian motion, so the purpose of the ratchet and pawl is to rectify the Brownian motion of these paddles. One may think that this could be used to do work (for example by using the axle to lift a weight), however this is not true. The problem is that the ratchet and pawl themselves will also be subject to random motion so they will sometimes allow the axle to turn anti-clockwise. To model the Feynman ratchet, we will need two degrees of freedom \cite{M.W.Jack2016}, this is beyond the scope of this project because we will only simulate systems with one degree of freedom.

\subsection{Landauer blowtorch}
The Landauer blowtorch scheme involves a temperature that varies in space \cite{Landauer1988}. As noted earlier, in order for the particles in Figure \ref{fig:Schematic} to get out of potential wells, they will need to acquire thermal energy from the environment. Landauer's idea was to assist these particles by heating the environment at the hills that the particles need to climb. From this Landauer notes that \cite{Landauer1988} ``The relative occupation of competing states of local stability is not determined solely by the characteristics of the locally favored states, but depends on the noise along the whole path connecting the competing states." This means that when we are modeling our system, we need to take non-constant temperatures into account.

\subsection{Thermal ratchet}
\cite{Pedro2014}

\subsection{Time dependent potentials}

\begin{figure}[tb]
	\centering
	\subfigure{%
		\includegraphics[width=0.45\columnwidth]{TimeDependentInit}
	}
\quad
	\subfigure{
		\includegraphics[width=0.45\columnwidth]{TimeDependentLater}
	}
	\subfigure{
		\includegraphics[width=0.45\columnwidth]{TimeDependentFinal}
	}
\caption{Schematic of particles diffusing in a time dependent potential, the blue line represents the potential and the black line represents the particle density. In (a) the particles are stuck in the first potential well, a certain time later the electric field is turned off and the particles diffuse freely as shown in part (b). When the potential is turned back on again, a large number of the particles will get stuck in the well to the left of the first one thus creating a net current to the left.}
\label{fig:TimeDependent}
\end{figure}

Both \cite{BlickleBechinger2011} \cite{JoelBader1999} are examples of experiments where a time dependent potential was used to do work on molecules. In \cite{JoelBader1999} DNA molecules are placed in a time dependent ratchet potential. When the electric field is on it creates a saw tooth shaped potential as shown in Figure \ref{fig:TimeDependent}, when the electric field is on, the molecules will go to the bottom of the wells and get stuck there. When the electric field is turned off however, the molecules will diffuse freely so that when the potential is turned back on at a later time, a large number of the molecules will then get trapped in the well to the left of where they were originally. Thus there is a net drift to the left.

\subsection{Tilted periodic potentials}
The tilted periodic potential (Figure \ref{fig:Schematic}) is of particular interest for this project because it can be used to model biological motors \cite{Leibler1993,Magnasco1994}. One way to model Brownian motors of this class is to think of a reaction coordinate $x$ that describes the conformation of a molecule in a chemical reaction. An example of this is the reaction ATP $\rightleftharpoons$ ADP + P, where ATP is adenosine tri-phosphate, ADP is adenosine di-phosphate and P is a lone phosphate molecule. This reaction coordinate is then coupled to a mechanical coordinate $y$ so that each time that a reaction takes place, the motor will move in some way. Since this is a chemical reaction, the free energy will be decreased as time moves forwards. So a system that has a potential that is periodic in $x$ is not sufficient to describe this situation, we will need to ``tilt" the potential by adding a forcing $f$. The value of $f$ will depend on the $\Delta G$ of the reaction (i.e. how far out of equilibrium the reaction is). It is shown in \cite{Magnasco1994} that this can be modeled by the two dimensional Smoluchowski equation. In this project we will only be modeling the one dimensional Somluchowski equation, so we will have to consider the case where $x$ and $y$ are tightly coupled. An example of tight coupling is the kinesin motor \cite{Leibler1993} that is used in cells to transport molecules. The kinesin motor is strongly bound to a track that it ``walks" along, on each step the motor will hydrolyze an ATP molecule using the reaction shown above. This reaction liberates about $12 k_B T$ Joules of energy that the motor uses to move forward. Kinesin motors are able to take many steps forward while taking few steps backwards all while falling off their track very infrequently \cite{BlockSM1990}.

\section{The Smoluchowski equation coupled to the environment} \label{Smoluchowski}

As we will see, the diffusion is increased by increasing the temperature and the drift is increased by increasing $f$. In some cases such as the Landauer blowtorch, the environment has a non uniform temperature held fixed by an external heat source \cite{Landauer1988}. With this in mind, we interpret Figure \ref{fig:Schematic} as follows: Brownian particles are subject to a given potential and are agitated by thermal noise, these agitations can give the particles the energy to move over barriers created by the potential. As one could imagine, these thermal interactions draw energy from the environment causing the temperature of the environment to change. Normally two simplifying assumptions are made at this point \cite{Reimann2001}, (i) that the thermal fluctuations created by the motor are very small compared to the thermal energy of the surrounding environment which is assumed to be effectively infinite, (ii) that when these temperature fluctuations occur, they diffuse away so rapidly that they do not need to be accounted for. In this project, we will question the second assumption in the case of Brownian motors. Assumption (ii) has also been questioned previously by Streater in the context of Brownian motion \cite{Streater1997, Streater1997a}. In these articles, Streater investigates Brownian motion from a microscopic view and then comes up with a mathematical model to describe Brownian particles that are thermally coupled to the environment, he then goes on to prove that the model is thermostatistically consistent in the sense that energy is conserved and that entropy increases. We will explore a similar set of equations in the context of Brownian motors and we will try to determine the length scales at which the thermal coupling is important.

The main part of this project will be trying to understand the behavior of the coupled partial differential equations given by:

\begin{eqnarray}
J(x, t) &=& \gamma^{-1} \frac{\partial}{\partial x} \left ( \frac{\partial V(x, t)}{\partial x} P(x, t) + k_B \frac{\partial}{\partial x} \left [T(x, t) P(x, t) \right] \right )  \\
\frac{\partial P(x, t)}{\partial t} &=& \frac{\partial J}{\partial x} \label{eqn:Smoluchowski} \\
\frac{\partial T(x, t)}{\partial t} &=& \frac{\partial}{\partial x} \left ( -\kappa q(x, t) + D \frac{\partial T(x, t)}{\partial x} \right ) \label{eqn:TemperatureEvolution}
\end{eqnarray}

Where
\begin{itemize}
\item{$P(x, t)$ is the probability density as a function of  reaction coordinate $x$ and time $t$}
\item{$J(x, t)$ is called the current}
\item{$\gamma$ is the friction coefficient}
\item{$V(x, t)$ is the potential for the motor}
\item{$k_B$ is the Boltzmann constant}
\item{$q(x, t) = \partial_x V(x, t) J(x, t)$ is the heat from the motor}
\item{$\kappa$ is the thermal conductivity}
\item{$D$ is the thermal diffusivity}
\end{itemize}

Equation (\ref{eqn:Smoluchowski}) is called the Smoluschowski equation \cite{KellerBustamante2000} and equation \ref{eqn:TemperatureEvolution} is the heat equation. These equations make our intuitive notions more precise, we see that the first term on the right hand side of the Smoluchowski equation (equation \ref{eqn:Smoluchowski}) is a drift term that is forced by our potential and that the second term contains a diffusion term that is scaled by our temperature. In fact, Figure \ref{fig:Schematic} was made by solving equation (\ref{eqn:Smoluchowski}) numerically. Likewise, equation (\ref{eqn:TemperatureEvolution}) also appeals to how our intuition of how the motor should effect its environment. The first term represents the heat flux being produced by the motor \cite{M.W.Jack2016}, while the second term represents the diffusion of temperature into the environment.


\chapter{Setting up the system}

\epigraph{The law that entropy always increases, holds, I think, the supreme position among the laws of Nature. If someone points out to you that your pet theory of the universe is in disagreement with Maxwell's equations — then so much the worse for Maxwell's equations. If it is found to be contradicted by observation — well, these experimentalists do bungle things sometimes. But if your theory is found to be against the second law of thermodynamics I can give you no hope; there is nothing for it but to collapse in deepest humiliation.}{\textbf{Arthur Eddington \\ -- The nature of the physical world (1928)}}


\section{The Smoluchowski equation} \label{Smoluchowski}
Here we would like to explain our model for Brownian dynamics with self-induced temperature gradients, if we know the temperature as a function of space and time, then we would use the Smoluchowski equation (equation \ref{eqn:smoluchowskiIntro}) to model the motion. One can interpret the relation between equation \ref{eqn:smoluchowskiIntro} and Figures \ref{fig:Schematic}, \ref{fig:landauersBlowtorch} as follows: Brownian particles are subject to a given potential and are agitated by thermal noise, these agitations can give the particles the energy to move over barriers created by the potential. As one could imagine, these thermal interactions draw energy from the environment as noted in Figure \ref{fig:brownianMotion}, one may imagine that these interactions could cause temperature of the environment to change. Normally two simplifying (not always explicit) assumptions are made at this point \cite{Reimann2001}, (i) that the temperature gradient created by the motor are very small compared to the thermal energy of the surrounding environment which is assumed to be effectively infinite, (ii) that when these temperature fluctuations occur, they diffuse away so rapidly that they do not need to be accounted for. In this project, we will question the second assumption and explore the consequences of not making this assumption.

First, we should derive the heat produced by the Brownian particle as it diffuse; this is done in Ref \cite{M.W.Jack2016}, here we will reproduce the derivation because it is very important. Notice that the potential energy of the motor is given by $U = \int P(x, \, t) V(x) dx$. Differentiating with respect to time, we get
\begin{eqnarray}
\frac{d U}{d t} &=& -\int \frac{\partial J(x, \, t)}{\partial x} V(x) dx \\
		        &=& -[J(x) V(x)] + \int J(x, \, t) \partial_x V(x, \, t) dx
\end{eqnarray}
The first term represents the work done on the particle and the integrand in the second term is the heat produced by the particle which we will denote with $\dot{q}(x, \, t)$. Now that we know the heat being produced by the particle, we are ready to quantify the self-induced temperature gradients. First of all, note that if there are no sources of heat, then the temperature will evolve according to the equation:
\begin{equation}
\frac{\partial T(x, \, t)}{\partial t} = 
\end{equation} 

%Assumption (ii) has also been questioned previously by Streater in the context of Brownian motion \cite{Streater1997, Streater1997a}. In these articles, Streater investigates Brownian motion from a microscopic view and then comes up with a mathematical model to describe Brownian particles that are thermally coupled to the environment, he then goes on to prove that the model is thermostatistically consistent in the sense that energy is conserved and that entropy increases. We will explore a similar set of equations in the context of Brownian motors and we will try to determine the length scales at which the thermal interaction is important.

The majority of the effort expended in this thesis will be in understanding the nature of two coupled partial differential equations, these take on the form:
\begin{eqnarray}
J(x, \, t) &=& -\gamma^{-1} \left ( \frac{\partial V(x)}{\partial x} P(x, \, t) + k_B T(x, \, t) \frac{\partial P(x, \, t)}{\partial x} \right ) \label{eqn:current} \\
\frac{\partial P(x, \, t)}{\partial t} &=& -\frac{\partial J}{\partial x} \label{eqn:Smoluchowski} \\
\frac{\partial T(x, \, t)}{\partial t} &=& -\kappa \dot{q}(x, \, t) + D \frac{\partial^2 T(x, 
, t)}{\partial x^2} \label{eqn:TemperatureEvolution}
\end{eqnarray}

Where
\begin{itemize}
\item{$P(x, \, t)$ is the probability density as a function of  reaction coordinate $x$ and time $t$}
\item{$J(x, \, t)$ is the probability current}
\item{$\gamma$ is the friction coefficient}
\item{$V(x)$ is the potential for the motor}
\item{$k_B$ is the Boltzmann constant, units $J K^{-1}$}
\item{$\dot{q}(x, \, t) = \partial_x V(x, \, t) J(x, \, t)$ is the heat from the motor \cite{M.W.Jack2016}}
\item{$\kappa$ is dictates how much the motor effects the environment, in ~\autoref{thermodynamics} we will show that $\kappa$ is one over the specific heat of the environment}
\item{$D$ is the thermal diffusivity and has units $m^2 s^{-1}$}
\end{itemize}
Equation \ref{eqn:Smoluchowski} is called the Smoluschowski equation which we noted earlier and equation \ref{eqn:TemperatureEvolution} is the heat equation. These equations make our intuitive notions more precise, we see that the first term on the right hand side of the Smoluchowski equation (equation \ref{eqn:Smoluchowski}) is a drift term that is forced by our potential and that the second term contains a diffusion term that is scaled by our temperature.

Therefore equation \ref{eqn:TemperatureEvolution} appeals to our intuition of how the Brownian particle should effect its environment. The first term represents the heat flux being produced by the motor, while the second term represents the diffusion of temperature into the environment.

One may notice that the system above is a non-linear system of equations, if one looks at the last term of \ref{eqn:current}, then one will find that the second spatial derivative of $P(x, \, t)$ is multiplied by the temperature, which due to equation \ref{eqn:TemperatureEvolution} is a function of $P(x, \, t)$. Therefore equation \ref{eqn:current} is a second order non-linear equation. Likewise if one looks at the first term of equation \ref{eqn:TemperatureEvolution}, then one will find that the temperature is multiplied by the first derivative of the probability distribution (which is a function of temperature), therefore equation \ref{eqn:TemperatureEvolution} is also non-linear. Because we are dealing with non-linear equations, we must take extra care when finding their solutions numerically, in ~\autoref{numerics} we will discuss the techniques used to solve the equations as well as analytical and convergence tests that we used to test the performance of these numerical techniques.

In fluid dynamics equation \ref{eqn:TemperatureEvolution} is called the heat equation, in our case we have neglected the term that represents the flow of the fluid that the Brownian particle is situated in. Including this flow would involve adding $u(x) \frac{\partial T}{\partial x}$ to the left hand side of equation \ref{eqn:TemperatureEvolution}, where $u(x)$ is velocity of the flow of the fluid. The role of fluid dynamics is discussed further in ~\autoref{fluidDynamics}, here we will note that the assumption that the fluid is stationary ($u(x) = 0$) means that the temperature gradients cannot be allowed to become too large because temperature gradients are known to create flows in fluids. Even though these flows may effect the validity of our model in real life we will neglect them for three main reasons. (i) They are not related to the physics that we would like to capture in this project since this is a project on statistical mechanics, not fluid dynamics. (ii) We would not like to depart too far from the explored literature, this is similar to reason (i) however here we would like to stress that our model is just an extension of previous models, extending the extension before understanding the extension could be premature. (iii) Including the flow into our equations of motion would increase the complexity of the system to beyond the scope of an honors project: although simply adding the fluid flow to equation \ref{eqn:TemperatureEvolution} may seem like an innocent operation, calculating the way in which the temperature gradients and the flow of the Brownian particle effect the fluid flow would mean adding a thrid equation to our already non-linear coupled system of equations.

This model includes a temperature that depends on $x$ and $t$, which has been explored in the literature \cite{Kramers1940}. Some authors write the current as \cite{Kramers1940,Kampen1988}
\begin{equation}
J(x, \, t) = -\gamma^{-1} \frac{\partial}{\partial x} \left (\frac{\partial V(x, \, t)}{\partial x} P(x, \, t) + k_B \frac{\partial}{\partial x}[T(x, \, t) P(x, \, t)] \right )  \label{eqn:diffCurrent}
\end{equation}
the difference being that they put the temperature inside the inner derivative. We will return to this difference in ~\autoref{thermodynamics}, where we will discuss the role that the current has in the generation of entropy.
Our model departs from previous work is that the temperature now depends on the evolution of the probability distribution as well. Coupled models of this type have been mentioned earlier by Streater \cite{Streater1997, Streater1997a,Streater2000,Streater1997b}, our work uses a similar model that was modified to suit our needs.

Intuitively, the first term in equation \ref{eqn:TemperatureEvolution} will create temperature gradients and the second term will cause these temperature gradients to flatten out. If the first term dominates, then the temperature gradients will grow to a very large size. In this case, we have to concede that our model will no longer be relevant, in particular our model does not include phase changes, so at the very least we require that the temperature gradients do not cause the fluid that the particle is in to freeze. This can be done if one makes sure that the second term is not significantly smaller than the first. In section ~\autoref{dimensionless} we will show that this puts a restriction on the characteristic energy of the potential. On the other hand, if the second term completely dominates, then any temperature gradients created by the Brownian particle will immediately diffuse away. In this regime, equation \ref{eqn:TemperatureEvolution} will not have an effect on the evolution of the system and the model will be reduced to the Smoluchowski equation by itself. The regime that we are interested in is the regime where the first and second terms are both important, in this regime the Brownian particle will have an effect on the environment but this effect will not cause the system to behave in a way that causes our model to break down. 

The other restriction that we must make to our model is that there is only a very small number of Brownian particles. This is because the temperature gradients can act as inter-particle interactions, to see this imagine that there are two particles moving about and creating temperature gradients. If one particle moves to a location where the other has created a temperature gradient, then it will be effected by the temperature gradient created by the other particle. Therefore, to model this, one would need to model the probability distributions of each particle separately and make sure that the interactions are accounted for. We stress this point because often people will use the Smoluchowski equation to describe the density of particles in a fluid. Interpreting $P(x, \, t)$ in this way in the presence of self induced temperature gradients would neglect the inter-particle interactions and would not be an accurate description. Therefore, we will be careful to make sure that $P(x, \, t)$ is used to model the probability distribution of a single colloidal particle, not many.  

Temperature gradients that are self induced in Brownian dynamics have been known for some time, a particularly interesting example is the Soret-Dufour effect \cite{Onsager1931,HortLinzLuecke1992,PiazzaGuarino2002,Santamaria-HolekGadomskiRubi2011}. In this effect, the diffusion of the concentration of a species of molecules causes a measurable temperature effect. Another interesting example is the reverse Landauer blowtorch \cite{DasDasBarikEtAl2015}, where a nett force displacing a Brownian particle ca create a temperature gradient. This however, is in contrast to the forward Landauer blowtorch where a temperature gradient creates a psuedo-force on the particles. We will explore the reverse Landauer blowtorch in ~\autoref{reverseLandauerBlowtorch} where we discuss how our model adds to previous work.

% THERMODYNAMICS
\section{System thermodynamics} \label{thermodynamics}
Equation \ref{eqn:Smoluchowski} and equation \ref{eqn:TemperatureEvolution} define equations of motion for our system, the system may be confined to a region $\Omega$ embedded in a larger environment which interacts with our system through the boundary conditions. In this section, we will show that our equations of motion obey the first and second laws of thermodynamics under certain conditions.

In the over-damped limit, the kinetic energy of the Brownian particle is negligible, therefore the energy of the system is completely described by the potential energy of the particle plus the thermal energy of the bath. The potential energy of the particle is $U_P = \int_{\Omega} V(x) P(x) dx$ and if the fluid that we are in is incompressible, then the thermal energy of the bath is $c_p \int_{\Omega} T(x) dx$, where $c_p$ is the specific heat capacity of the environment, to see that this is the thermal energy consider the units of the quantity $c_p T(x, \, t)$. These are $J K^{-1} K$, so $c_p T(x, \, t) dx$ is the heat content of an infinitesimal element $dx$. With this we have:
\begin{equation}
E(t) = \int_{\Omega} V(x)P(x, t) dx + c_p \int_{\Omega} T(x, t) dx
\end{equation}
By using the Smoluchowski equation and the heat equation, we can differentiate with respect to time to get:
\begin{align}
\frac{d E}{d t} & = \int_{\Omega} V(x) \frac{\partial P}{\partial t} dx + c_p \int_{\Omega} \frac{\partial T}{\partial t} dx \\
 & = -\int_{\Omega} V(x) \frac{\partial J}{\partial x} + c_p \int_{\Omega} -\kappa J(x) \frac{\partial V}{\partial x} + D \frac{\partial^2 T}{\partial x^2} dx \\
 & = [V(x)J(x)]_{\partial \Omega}+ \int_{\Omega} \frac{\partial V}{\partial x} J(x) dx - \kappa c_p \int_{\Omega} \frac{\partial V}{\partial x} J(x) dx + c_p D \left [\frac{\partial T}{\partial x} \right]_{\partial \Omega}.
\end{align}
Notice that if the middle two terms cancel, then the change in energy is equal to the flow of energy through the boundaries. Thus, the first law of thermodynamics requires that $\kappa = \frac{1}{c_p}$, physically we can understand this by looking at the first term of equation \ref{eqn:TemperatureEvolution}. When the heat capacity is small, $\kappa$ becomes large and in this case, even a small amount of heat being produced by the Brownian particle will have a large effect on the evolution of the temperature. Conversely, if the heat capacity is large, then the heat produced by the Brownian particle will have a very small effect on the temperature. We therefore expect that equation \ref{eqn:TemperatureEvolution} can be neglected in the case where the environment has a very large heat capacity compared to the heat being produced by the Brownian particle.

As for the entropy, we have \cite{Streater1997a}
\begin{equation}
S(t) = -k_B \int_{\Omega} P(x, t) \log(P(x, \, t)) dx + c_p \int_{\Omega} \log(T(x, \, t))dx.
\end{equation}
The first term is the Gibbs entropy in the continuous case \cite{Jaynes1965} and the second term is the entropy of an incompressible fluid \cite{CengelBoles1994}. Using the entropy of an incompressible fluid restricts the sytems that can be modelled by our equations of motion, in particular we will not be able to treat Brownian particles suspended in a gas since a gas is compressible by defintion. The courageous reader may want to use our model in the case of a compressible fluid, however in this case we can no longer guarantee this reader that entropy will increase for them. The goal of this project is not to create a system that models many different situations in nature, but rather to create a system that is self-consistent.

Differentiating the entropy with respect to time, we get

\begin{align}
\frac{d S}{d t} =  k_B \int_{\Omega} \frac{\partial J}{\partial x} + \frac{\partial J}{\partial x} \log P \ dx + c_p \int_{\Omega} \frac{1}{T} \left(-\kappa J \partial_x V + c_p D \frac{\partial^2 T}{\partial x^2} \right) dx \\
                     = k_B \left ( [J]_{\partial \Omega} + [J \log P]_{\partial \Omega} - \int_{\Omega} \frac{J}{P} \frac{\partial P}{\partial x} dx \right) - \int_{\Omega} \frac{J}{T} \frac{\partial V}{\partial x} + c_p D \int_{\Omega} \frac{1}{T} \frac{\partial^2 T}{\partial x^2} dx
\end{align}
Denote the boundary terms as $B(t) = k_B( [J \log P]_{\partial \Omega} + \left[\frac{\partial J}{\partial x} \right]_{\partial \Omega} ) $, we can write:

\begin{eqnarray}
\frac{d S}{d t} & = & - \int_{\Omega} k_B \frac{J}{P} \frac{\partial P}{\partial x} + \frac{J}{T} \frac{\partial V}{\partial x} dx +  c_p D \int_{\Omega} \frac{1}{T} \frac{\partial^2 T}{\partial x^2} dx + B(t) \\
                    & = & \gamma \int_{\Omega} \frac{J^2}{T P} dx + c_p D \int_{\Omega} \frac{1}{T} \frac{\partial^2 T}{\partial x^2} dx + B(t), \label{eqn:entropyGen}
\end{eqnarray}
where in the second equality we used the fact that $J = -\gamma^{-1} ( P \frac{\partial V}{\partial x} + k_B T \frac{\partial P}{\partial x})$.

By noticing that
\begin{equation}
\frac{\partial}{\partial x} \left(\frac{1}{T} \frac{\partial T}{\partial x} \right) = -\frac{1}{T^2} \left(\frac{\partial T}{\partial x} \right)^2 + \frac{1}{T} \frac{\partial^2 T}{\partial x^2},
\end{equation}
we can rewrite the second term of equation \ref{eqn:entropyGen} as:
\begin{eqnarray}
c_p D \int_{\Omega} \frac{1}{T} \frac{\partial^2 T}{\partial x^2} dx &=& c_p D \int_{\Omega} \frac{\partial}{\partial x} \left(\frac{1}{T} \frac{\partial T}{\partial x} \right) + \frac{1}{T^2} \left(\frac{\partial T}{\partial x} \right)^2 dx \\
 &=& c_p D \int_{\Omega} \frac{1}{T^2} \left(\frac{\partial T}{\partial x} \right)^2 dx + c_p D\left[\frac{1}{T} \frac{\partial T}{\partial x} \right]_{\partial \Omega}
\end{eqnarray}
Define
\begin{eqnarray}
\dot{S}_{gen} &\equiv& \int_{\Omega} \gamma \frac{J^2}{T P} + c_p D \frac{1}{T^2} \left(\frac{\partial T}{\partial x} \right)^2 dx \label{eqn:entropyGen} \\
\text{and} \nonumber \\
B(t) &\equiv& k_B \left( [J \log P]_{\partial \Omega} + \left[\frac{\partial J}{\partial x}\right]_{\partial \Omega} \right) + c_p D\left[\frac{1}{T} \frac{\partial T}{\partial x} \right]_{\partial \Omega}
\end{eqnarray}
so the change in entropy can be written as:
\begin{equation}
\frac{d S}{d t} = \dot{S}_{gen} + B(t)
\end{equation}
We notice that the change in entropy is equal to a positive number $\dot{S}_{gen}$ plus the entropy flowing through the boundaries, this is precisely the second law of thermodynamics. Furthermore, the generated entropy can be split into two terms, the first term is to be interpreted as the entropy generated by the Brownian particle; since the motion of the particle is random, as the Brownian particle diffuses we will become less certain of its position. The second term is the entropy generated by the diffusion of temperature gradients; if one studies the solutions to the heat equation, then one will find that any spikes in the temperature will flatten out and eventually the temperature will be uniform. This flattening out of the temperature is irreversible and is therefore associated with an increase in entropy. When one discusses heat engines implemented on the molecular scale, one finds that the Carnot efficiency can only be attained when the entropy generation is zero. Simply by reading equation \ref{eqn:entropyGen}, we can see that this occurs when the current is zero and the temperature is flat. In ~\autoref{SteadyState} we will show that this is not possible for a Brownian particle in a tilted potential. 

As mentioned earlier, some authors write the current in the form of equation \ref{eqn:diffCurrent}, however if we use this definition of $J$, then $\dot{S}_{gen}$ is not necessarily positive, in fact with this term in the equation $\dot{S}_{gen}$ becomes:
\begin{equation}
\dot{S}_{gen} \equiv \int_{\Omega} \gamma \left(\frac{J^2}{T P} - \frac{J}{T} \frac{\partial T}{\partial x} \right) + c_p D \frac{1}{T^2} \left(\frac{\partial T}{\partial x} \right)^2 dx 
\end{equation}
With this version of the current, we can get local decreases in entropy, therefore we do not use this version of the current.

Our version of the current is supported by Streater \cite{Streater1997, Streater1997a,Streater2000,Streater1997b} and will reduce to the other version in the case of constant temperature.

% DIMENSIONLESS
\section{Making the equations dimensionless}  \label{dimensionless}

Upon viewing equations \ref{eqn:Smoluchowski} and \ref{eqn:TemperatureEvolution}, we see that there is a large number of constants that are set by the properties of the Brownian particle that we are modeling. We would like to reduce the number of variables for two reasons (i) by reducing the number of variables we will hopefully gain a more concise physical description of the system (ii) having a small number of free variables is very convenient for creating a program to approximate the equations numerically, dimensionless equations tend to be less prone to numerical error because they avoid cases where small numbers are compared to large ones in a floating point system.
Here we will non-dimensionalize the equations, to do this, introduce $\hat{x} = \frac{x}{L}$, where $L$ is the length scale of the system, the Smoluchowski equation now becomes

\begin{equation}
\frac{\partial P}{\partial t} = \gamma^{-1}\frac{1}{L^2} \frac{\partial}{\partial \hat{x}} \left (P \frac{\partial V}{\partial \hat{x}} + k_B T \frac{\partial P}{\partial \hat{x}} \right )
\end{equation}
Now let $E_0$ be the characteristic energy of the system, for a chemical reaction, $E_0$ will be the $\Delta G$ of the reaction and for a periodic potential, $E_0$ will be amplitude of the oscillations. Now we introduce the dimensionless potential and the dimensionless temperature as $\hat{V}(x) = \frac{V(x)}{E_0}$ and $\hat{T}(x) = \frac{k_B T(x)}{E_0}$ respectively. Now the Smoluchowski equation becomes

\begin{equation}
\frac{\partial P}{\partial t} = \frac{E_0}{\gamma L^2} \frac{\partial}{\partial \hat{x}} \left (P \frac{\partial \hat{V}}{\partial \hat{x}} + \hat{T} \frac{\partial P}{\partial \hat{x}} \right )
\end{equation}
Let $\hat{t} = \frac{E_0}{\gamma L^2} t$ and $\hat{P} = L P $, with this we have:

\begin{equation}
\frac{\partial \hat{P}}{\partial \hat{t}} = \frac{\partial}{\partial \hat{x}} \left (\hat{P} \frac{\partial \hat{V}}{\partial \hat{x}} + \hat{T}  \frac{\partial \hat{P}}{\partial \hat{x}} \right ) \label{eqn:dimensionlessSmoluchowski}
\end{equation}
Now we define $\hat{J} = \frac{\gamma L^2}{E_0} J $, applying these definitions to the equation for the temperature evolution, we find that:

\begin{equation}
\frac{E_0^2}{\gamma k_B L^2} \frac{\partial \hat{T}}{\partial \hat{t}} = -\kappa \frac{E_0}{\gamma L^2}\hat{J}(\hat{x}) \frac{E_0}{L} \frac{\partial \hat{V}}{\partial \hat{x}} + \frac{D E_0}{k_B L^2} \frac{\partial^2 \hat{T}}{\partial \hat{x}^2}
\end{equation}
Now let 
\begin{equation}
\alpha = \frac{\kappa k_B}{L} \quad \text{or} \quad \alpha = \frac{k_B}{c_p L} \\
\end{equation}
and
\begin{equation}
\beta = \frac{D \gamma}{E_0},
\end{equation}
then we have
\begin{equation}
\frac{\partial \hat{T}}{\partial \tau} = -\alpha \hat{J}(\hat{x}) \frac{\partial \hat{V}}{\partial \hat{x}} + \beta \frac{\partial^2 \hat{T}}{\partial \hat{x}^2} \label{eqn:dimensionlessHeat}
\end{equation}

As for the energy of the system, the dimensioned version is:

\begin{equation}
E(t) = \int P(x) V(x) dx + \frac{1}{\kappa} \int T(x) dx
\end{equation}
Let $\hat{E}(t) = \frac{E(t)}{E_0}$, we have:

\begin{eqnarray}
\hat{E}(t) & = & \int \hat{P} \hat{V} d \hat{x} + \frac{L}{k_B \kappa} \int \hat{T} d\hat{x} \\
              & = & \int \hat{P} \hat{V} d \hat{x} + \frac{1}{\alpha} \int \hat{T} d\hat{x}
\end{eqnarray}

So our system depends on the parameters $\alpha$ and $\beta$ as well as the shape of the potential. Physically, $\alpha$ represents how much the particle interacts with the environment thermally and $\beta$ represents how quickly the temperature gradients diffuse. As the length scale of the system increases, $\alpha$ will decrease, therefore the temperature gradients become negligible for large systems. We interpret this as meaning that the temperature gradients are much smaller than the system and can diffuse away quickly. As the characteristic energy of the system increases, $\beta$ will decrease meaning that the diffusion of temperature will be negligible. We interpret this as meaning that for high energy systems, temperature gradients are created much faster than they can be diffused away. 


%In order to neglect the first term of \ref{eqn:dimensionlessHeat}, we need to have $\alpha \hat{J} \hat{V} \ll \frac{\partial \hat{T}}{\partial \tau}$, this means that  we must have:
%\begin{equation}
%\frac{k_B}{c_P L} \frac{\gamma L^2}{E_0} J(x, \, t) \frac{L}{E_0} \partial_x V(x, \, t) \ll \left (\frac{E_0}{\gamma L^2} \right)^{-1} \frac{E_0}{k_B} \frac{\partial T(x, \, t)}{\partial t}
%\end{equation}
%So in terms of the physical properties of the physical parameters of the environment,
%\begin{equation}
%\frac{1}{c_p} \left (\frac{k_B}{E_0} \right)^2 \ll (J(x, \, t) \partial_x V(x, \, t) )^{-1} \frac{\partial T(x, \, t)}{\partial t} \label{eqn:neglectAlpha}
%\end{equation}
%We can understand this physically as follows: $E_0$ represents the typical barrier height that the Brownian particle will encounter in it's energy landscape. When these barriers are large compared to the energy of the particle, the particle will essentially be trapped and the temperature gradients that it produces will be negligible. Also, when the specific heat of the environment is large, any heat that the particle draws from the bat will have a negligible effect on the actual temperature of the environment.

%Meanwhile, in order for us to neglect the second term, we need to have:
%\begin{equation}
%\frac{D \gamma}{E_0} L^2 \frac{\partial^2 \hat{T}(x, \, t)}{\partial x^2} \ll \left (\frac{E_0}{\gamma L^2} \right)^{-1} \frac{\partial \hat{T}(x, \, t)}{\partial t}
%\end{equation}
%This means that
%\begin{equation}
%D \ll 1  \label{eqn:neglectBeta}
%\end{equation}
%or in other words, the length scale of the system is much larger than the thermal diffusivity. Physically this means that if there are any initial thermal gradients, then they will remain in place over time because the thermal diffusion is too slow to dissipate them. Interestingly this regime does not necessarily mean that temperature gradients should be neglected entirely, this is because $\beta$ is independent of $\alpha$. If we meet the condition in equation \ref{eqn:neglectBeta}, but at the same time we do not meet the condition in equation \ref{eqn:neglectAlpha}, then the Brownian particle will create temperature gradients that will not diffuse away rapidly. In these cases, the small thermal gradients may diverge to arbitrary sizes as time goes to infinity. Our model is only valid for small deviations in temperature, in particular our model does not consider any phase changes so the temperature gradients must be constrained so that they do not cause the fluid to freeze. We will therefore avoid these situations in our numerical exploration of the system.

%There is one more condition that can be met under which the system effectively reduces to the Smoluchowski equation, imagine that we have a system that begins with a constant temperature in space. As time goes on, the term involving $\alpha$ induces temperature gradients and the term involving $\beta$ will cause these temperature gradients to flatten out. If the time scale of the thermal diffusion is much faster than the time scale of the production of thermal gradients, then the temperature gradients will all diffuse faster than they can be created, we can describe this mathematically by:
%\begin{equation}
%\alpha \hat{J} \partial_{\bar{x}} \hat{V} \ll \beta \frac{\partial \hat{T}(\bar{x}, \, t)}{\partial \bar{x}}
%\end{equation}
%we can reduce this down to the following condition:
%\begin{equation}
%\frac{k_B}{c_p E_0 D} \ll 1
%\end{equation}
%If this condition is true, then any temperature gradients that arise will diffuse away very quickly meaning that our model will reduce to the Smoluchowski equation.

%So far we have shown the conditions that allow us to either neglect specific terms in equation \ref{eqn:dimensionlessHeat}, but now we would like to discuss those conditions under which we need to use this equation in order to properly model a system. Currently we can say when certain terms can be neglected, but we would like to be able to determine when they are important. Specifically, we would like to know the conditions when the both terms are non-negligible and when the temperature gradients do not rapidly diffuse away. This means that all of the above conditions must be violated, we will write this as
%\begin{eqnarray}
%\frac{E_0}{L} &\sim& \frac{k_B}{c_p} \\
%L &\sim& D \\
%E_0 &\sim& \frac{k_B}{c_p D}
%\end{eqnarray}
%this is an upper bound for both $E_0$ and $L$ as well as an upper bound for their ratio.
%
%At standard temperature and pressure ($T = 298.15 K$ and pressure $P_{atm} = 101.3 kPa$), water has the following properties:
%\begin{itemize}
%\item{Specific heat: 4186 $J K^{-1}$}
%\item{Thermal diffusivity: $1.43 \times 10^{-5} m^2 s^{-1}$}
%\end{itemize}
%we also note that as per usual the Boltzmann constant is given by $k_B = 1.381 \times 10^{-23} J K^{-1}$. Typical chemical reactions that occur in biological molecular motors have energy on the order of $k_B T \approx 3.451 \times 10^{-22} J$ {\color{red} [citations]}. The upper bound for our energy is $E_{\text{upper}} \sim \frac{k_B}{c_p D} \approx \frac{1.381 \times 10^{-23}}{4186 \times 1.43 \times 10^{-5}} \approx 2.31 \time 10^{-22} J$. This means that molecular motors operate at the energy scale where equation \ref{eqn:dimensionlessHeat} is important, as for length scale 

\section{Model assumptions}
We are now in a position that we may discuss the physics of the equations of motion in some detail. In particular we would like to discuss the assumptions that have to be made in order for the equations to be valid. These assumptions have been explained sporadically throughout the text and we would like to compile them into a short that list that the reader should keep in mind throughout the remainder of the thesis.
\begin{enumerate}
\item{The degrees of freedom of the bath operate at a time scale much faster than that of the Brownian particle itself}
\item{The kinetic energy of the particle is negligible (over-damped limit)}
\item{There is only one Brownian particle in our model}
\item{Heat produced by chemical reactions in the power stroke of the motor is negligible}
\item{The fluid that the particle is in is incompressible}
\item{The fluid that the particle is in is steady}
\item{The temperature gradients are small enough that they do not induce a phase change}
\end{enumerate}
These assumptions have been discussed in the text, we can take any of these phenomena into account by making changes to our equations of motion, however these small changes to the equation of motion can cause the system to become very complex. Having said that, some of the assumptions made turn out to be more reasonable than one may have thought on first sight. Assumptions 1 and 2 are assumptions that are made when deriving the Smoluchowski equation from the Langevin equation, these assumptions have been discussed at length and turn out to be very reasonable \cite{Reimann2001,Gardiner2009,Einstein1905,KellerBustamante2000,Kramers1940}. Meanwhile, assumption 3 means that we cannot talk about the very interesting cooperative effects that occur in motors such as myosin, however this assumption is fine for motors such as kinesin or for colloidal particles in a potential. Assumption 4 {\color{red} talk to Michael Jack}. Assumption 5 is perfectly fine in water which is the fluid of choice for all biological molecular motors. Assumption 6 is fine as long as we do not have significantly large temperature gradients, finally assumption 7 just puts a restriction on the parameters that we choose in our simulations. As far as we are aware, there are no experimentally realized examples of Brownian particles creating a phase change in the fluid that they are suspended in.







\chapter{Solving the system}

\section{Steady state solution} \label{SteadyState}
In order to see how the equations of motion behave with time, we have to resort to numerical methods (see section \ref{numerics}). However we note that for a given potential there will be a stationary solution that we will refer to as the ``steady state'', given periodic boundary conditions, we will derive an analytical form for this steady state.

In the steady state, we have:

\begin{eqnarray}
\frac{\partial P(x, t)}{\partial t} &=&  0 \ = \ \frac{\partial J}{\partial x} \label{eqn:SmoluchowskiSteady} \\
\frac{\partial T(x, t)}{\partial t} &=& 0 \ = \ -\kappa q(x,t) + \frac{\partial}{\partial x} \left ( D \frac{\partial T(x, t)}{\partial x} \right ) \label{eqn:TemperatureSteady}
\end{eqnarray}

Suppose that we have the following boundary conditions:
\begin{align}
P(x = 0) & = P(x = L) \\
J(x = 0) & = J(x = L) \\
\left. \frac{\partial T}{\partial x} \right \rvert_{x = 0} & = 0 = \left. \frac{\partial T}{\partial x} \right \rvert_{x = L} \label{eqn:temperatureBoundary}
\end{align}
where $L$ is the length scale of the system. Physically, these conditions say that the nett current flowing out of the boundaries is zero and that no heat escapes from the system, thus the energy of the sytem is conserved. Section 5.2 of \cite{Gardiner2009} gives the steady state current as:

\begin{equation}
J_s = \left [\frac{2 k_B T(L)}{\psi(L)} - \frac{2 k_B T(0)}{\psi(0)}  \right] P_s(0) \left [\int_0^L dx'/\psi(x') \right]^{-1}
\label{eqn:SteadyCurrent}
\end{equation}

with $\psi(x) \equiv \exp[-\int_0^x dx' \frac{\partial_x V(x')}{2 k_B T(x')}]$. Meanwhile, the density is:

\begin{equation}
P_s(x) = P_s(0) \left [\frac{\int_0^x \frac{dx'}{\psi(x')} \frac{T(L)}{\psi(L)} + \int_x^L \frac{dx'}{\psi(x')} \frac{T(0)}{\psi(0)} }{\frac{T(x)}{\psi(x)} \int_0^L \frac{dx'}{\psi(x')} } \right]
\label{eqn:SteadyDensity}
\end{equation}
In this case, $J_s$ is a constant and $P_s(0)$ is also a constant. Assuming that we know these constants it is now possible to find the steady state temperature. We have:

\begin{equation}
\frac{\partial T}{\partial t} = 0 = -\kappa J_s \partial_x V + D \frac{\partial^2 T}{\partial x^2} \label{eqn:steadyTemperatureODE}
\end{equation}

In one dimension, \ref{eqn:steadyTemperatureODE} can be written as an ordinary differential equation of the form
\begin{equation}
T''(x) = \frac{\kappa J_s}{D} V'(x)
\end{equation}

We can solve this equation by integrating both sides twice to give:

\begin{equation}
T(x) = \frac{\kappa J_s}{D} \int_0^x V(x') dx' + \xi x + d \label{eqn:steadyTemperature}
\end{equation}
for unkown constants $\xi$ and $d$. By applying the boundary condition \ref{eqn:temperatureBoundary}, we find that
\begin{align}
T'(0) & = 0 = \frac{\kappa J_s}{D} V(0) + \xi \\
T'(L) & = 0 = \frac{\kappa J_s}{D} V(L) + \xi
\end{align}
This implies that either $J_s = 0$ or $V(0) = V(L)$, meaning that the coupled system does not admit a steady state solution for a tilted potential. Later, we will see that it is possible to have a steady state in higher dimensions where the flow of heat from the environment can dissipate the heat produced by the Brownian particle. By recalling that $E = \int_0^L P(x') V(x') dx' + c \int_0^L T(x') dx'$ we are able to find an expression for $d$. First, we will integrate the temperature from $0$ to $L$
\begin{equation}
\int_0^L T(x') dx' = \frac{\kappa J_s}{D} \int_0^L dx' \int_0^{x'} V(x'') dx'' + \xi \frac{L^2}{2} + d L
\end{equation}
Therefore,
\begin{equation}
	d = \frac{1}{L} \left(c E - c\int_0^L P(x') V(x') dx' - \frac{\kappa J_s}{D} \int_0^L dx' \int_0^{x'} V(x'') d x'' + \xi \frac{L^2}{2} \right)
\end{equation}

It would seem that one should be able to calculate the steady state current and density directly from the equations shown above. However, we notice that the constants $J_s$ and $P_s(0)$ have to satisfy equations (\ref{eqn:SteadyCurrent}), (\ref{eqn:SteadyDensity}) and (\ref{eqn:steadyTemperature}) while also satisfying the normalization condition $\int_0^L P(x) dx = 1$. To do this we define an objective function given by

\begin{equation}
obj(J_s, P_s(0)) = \left (J_s - \left [\frac{2 k_B T(L)}{\psi(L)} - \frac{2 k_B T(0)}{\psi(0)}  \right] P_s(0) \left [\int_0^L dx'/\psi(x') \right]^{-1} \right)^2  \label{eqn:Objective}
\end{equation}

And we minimize this objective function with respect to $J_s$ and $P_s(0)$ under the constraint $\int_0^L P(x) dx = 1$. Another way to do this is to guess a steady state density and temperature and use finite differencing to simulate forward in time until the transients die out.

%----------------------------------------------------------------------------------------
%	NUMERICS
%----------------------------------------------------------------------------------------
\section{Finite differences}  \label{numerics}
The one dimensional equation can be solved on a discrete grid by using the finite differences method, the main idea behind this strategy is to approximate derivatives with equations of the form:

\begin{equation}
\frac{d f}{d x} \approx \frac{f(x - h) - f(x + h)}{2h}
\end{equation}

for some small $h$, likewise the second derivative of a function is approximated with
\begin{equation}
\frac{d^2 f}{d x^2} = \frac{f(x - h) - 2f(x) + f(x + h)}{h^2}
\end{equation}
In our simulations, we will use the Crank Nicolson scheme to solve the equations. From now on, we will use the notation that $F(j \Delta x, n \Delta t) = F_j^n$, the key equation for the Crank Nicolson scheme is:
\begin{equation}
\frac{P_j^{n+1} - P_j^n}{\Delta t} = \frac{1}{2}(F_j^{n+1} + F_j^n)
\end{equation}
where $F$ represents the right hand side of the equation that we are doing finite differences on. By applying finite differences to the dimensionless Smoluchowski equation (eq \ref{eqn:dimensionlessSmoluchowski}), we find that:
\begin{equation}
F_j^{i} = \frac{P^i_{j+1} \partial V^i_{j+1} - P^i_{j-1} \partial V^i_{j-1}}{2 \Delta x} + \frac{T^i_{j+1} - T^i_{j-1}}{2 \Delta x} \frac{P^i_{j+1} - P^i_{j-1}}{2 \Delta x} + T^i_j \frac{}{} \frac{P^i_{j+1}- 2P^i_j + P^i_{j-1}}{\Delta x^2} 
\end{equation}
where we omitted the hats for notational convenience.
We make the following definitions:
\begin{align*}
a_j^{n+1} &= \frac{-2 T^{n+1}_j}{\Delta x^2} \\
b_j^{n+1} &=  \frac{\partial_x V^{n+1}_{j+1}}{2\Delta x} + \frac{T^{n+1}_{j+1} - T^{n+1}_{j-1}}{4 \Delta x^2} \\
c_j^{n+1} &= -\frac{\partial_x V^{n+1}_{j-1}}{2\Delta x}  - \frac{T^{n+1}_{j+1} - T^{n+1}_{j-1}}{4 \Delta x^2} \\ 
a_j^{n} &= \frac{-2 T^{n}_j}{\Delta x^2} \\
b_j^{n} &=  \frac{\partial_x V^{n}_{j+1}}{2\Delta x} + \frac{T^{n}_{j+1} - T^{n}_{j-1}}{4 \Delta x^2} \\
c_j^{n} &= -\frac{\partial_x V^{n}_{j-1}}{2\Delta x}  - \frac{T^{n}_{j+1} - T^{n}_{j-1}}{4 \Delta x^2} \numberthis
\end{align*}
With these definitions, the Crank Nicolson scheme can be written down as follows:
\begin{multline}
 -\frac{\Delta t}{2}a_j^{n+1}P_{j-1}^{n+1} + \left (1 - \frac{\Delta t}{2}b_j^{n+1} \right) P_j^{n+1} - \frac{\Delta t}{2} c_j^{n+1} P_{j+1}^{n+1} = a_j^n P_{j-1}^{n}
+ \left (1 + \frac{\Delta t}{2}b_j^n \right) P_j^{n}  + \frac{\Delta t}{2} c_j^n P_{j+1}^{n}
\end{multline}
This equation can be written in matrix form by defining the following matrices:
\begin{equation}
A =
\begin{bmatrix}
	a_0^{n+1} & b_1^{n+1} & 0                & 0                & 0        & \dots                                 & 0        \\
	c_0^{n+1} & a_1^{n+1} & b_2^{n+1} & 0                & 0        & \dots                                 & 0        \\
	0                & c_1^{n+1} & a_2^{n+1} & b_3^{n+1}   & 0        & \dots                                 & 0        \\
			     &                   &                   &                  &           &                                         &           \\
	\vdots         & \vdots         & \ddots         & \ddots         & \ddots & \vdots                              & \vdots \\
			     &                   &                   &                  &           &                                         &           \\
			     &                   &                   &                  &           &                                         &           \\
	0                &                   & \dots           &                   &  c_{J-2}^{n+1} & a_{J-1}^{n+1}  & b_J^{n+1} \\
	0                &                   & \dots           &                   &                         &  c_{J-1}^{n+1} & a_J^{n+1}
\end{bmatrix}
,\quad P^{n+1} =
\begin{bmatrix}
P_0^{n+1}       \\
P_1^{n+1}	     \\
                        \\
                        \\
\vdots               \\
                        \\
                        \\
P_{J-1}^{n+1} \\
P_J^{n+1}
\end{bmatrix}
\end{equation}

\begin{equation}
B =
\begin{bmatrix}
	a_0^{n} & b_1^{n}     & 0                 & 0          & 0                    & \dots            & 0        \\
	c_0^{n} & a_1^{n}     & b_2^{n}      & 0          & 0                    & \dots            & 0        \\
	0                & c_1^{n} & a_2^{n}      & b_3{n} & 0                    & \dots            & 0        \\
			     &               &                   &             &                      &                     &           \\
	\vdots         & \vdots     & \ddots         & \ddots   & \ddots            & \vdots           & \vdots \\
			     &               &                   &             &                      &                     &           \\
			     &               &                   &             &                      &                     &           \\
	0                &               & \dots           &             &  c_{J-2}^{n} & a_{J-1}^{n}  & b_J^{n} \\
	0                &               & \dots           &             &                     &  c_{J-1}^{n} & a_J^{n}
\end{bmatrix}
,\quad P^{n} =
\begin{bmatrix}
P_0^{n}       \\
P_1^n          \\
                    \\
                    \\
\vdots           \\
                    \\
                    \\
P_{J-1}^{n} \\
P_J^{n}
\end{bmatrix}
\end{equation}
With these matrices the equation now becomes,
\begin{equation}
\left (\mathbb{1} - \frac{\Delta t}{2}A \right) \cdot P^{n+1} = \left (\mathbb{1} + \frac{\Delta t}{2}B \right) \cdot P^n
\end{equation}
we interpret this equation as saying that half of a backwards Euler step acting on $P^{n+1}$ is equal to half of a forward Euler step acting on $P^n$. We write the equation to step $P$ forward one time step as
\begin{equation}
P^{n+1} = \frac{\mathbb{1} + \frac{\Delta t}{2}B}{\mathbb{1} - \frac{\Delta t}{2}A} P^n
\end{equation}
Each time that we step forward using this equation we will be out by a factor, this means that at each step we will need to renormalize using the equation $\int P(x) dx = 1$.

Likewise, we can apply the Crank Nicolson scheme to the heat equation, (eq \ref{eqn:dimensionlessHeat}), by looking at the right hand side of this equation, we find that

\begin{equation}
F_j^i = -\alpha \left (P_j^i (\partial_x V_j^i)^2 + \frac{T_{j+1}^i - T_{j-1}^i }{2 \Delta x}\partial_x V_j^i  \right) + \beta \frac{T_{j+1}^i - T_j^i + T_{j-1}^i}{\Delta x^2}
\end{equation}

Just like the discretized Smoluchowski equation, these equations can be written in matrix form. The temperature is normalized by assuming that the energy remains fixed, this will be true as long as no heat or current flows through the boundaries, i.e. $J(x=a) = 0 = J(x = b)$ and $\frac{\partial T}{\partial x} \rvert_a = 0 = \frac{\partial T}{\partial x} \rvert_b$. In this case, the energy is constant and is given by $E = \int P(x) V(x) dx + c_p \int T(x) dx$, so each time that we step the temperature forward, we have to calculate the potential and thermal energy and then scale the temperature so that the total energy remains fixed.

Fortunately the matrices that we are dealing with are very sparse, so the program used to solve these equations can save on memory by calling sparse matrix libraries. 

\section{Testing the numerics}

The idea behind finite differences is that as the discretization size goes to zero, the numerical approximation should converge on the correct analytical solution. We will compare our numerics with some known analytical results as well as performing convergence tests.
\subsection{A comparison with analytical results}

The Smoluchowski equation has a steady state probability density that takes on the form
\begin{equation}
P_{ss}(x) = N \exp{\left(-\int_a^x \frac{V'(x')}{T(x')} dx' \right)} \label{eqn:analSteadyState}
\end{equation}

Figure \ref{fig:smoluchowskiCompare} shows the a simulation where we began the system in the analytically calculated steady state, we then used finite differences to step forward 50,000 steps with $\Delta t = 10^{-4}$. After this simulation, we only found a minimal divergence from the steady state.
\begin{figure}
	\begin{subfigure}{0.49\textwidth}
	\includegraphics[width=\columnwidth]{smoluchowskiAnalytic}
	\end{subfigure}
	\begin{subfigure}{0.49\textwidth}
	\includegraphics[width=\columnwidth]{smoluchowskiAnalyticNormDiff}
	\end{subfigure}
	\caption{Finite differences in the steady state: The steady state for the system was calculated using equation ~\ref{eqn:analSteadyState}, we start the system off in this state and then simulate the system forward 50,000 steps with $\Delta t = 10^{-4}$ and $\Delta x = 0.02$. (a) shows the analytical steady state with the state after the simulation, (b) shows the absolute difference between these two vectors. Even after 50,000 steps the system has not deviated from the analytical steady state significantly. \label{fig:smoluchowskiCompare}}
\end{figure}

The heat equation can be solved using a Fourier series technique
\begin{equation}
\frac{\partial T}{\partial t} = \beta \frac{\partial^2 T}{\partial x^2}
\end{equation}
the initial condition for the temperature will be denoted by,
\begin{equation}
T(x, 0) = f(x)
\end{equation}

Lets say that the boundaries are at $x = \pm \infty$ and that the derivative of the temperature is zeros at the boudaries. The solution to the heat equation is given by:

\begin{equation}
T(x, t) = \sum_{n=1}^\infty D_n \sin \left(\frac{n \pi x}{L} \right) \exp\left(\frac{-n^2 \pi^2 \beta t}{L^2}\right)  \label{eqn:analTemperature}
\end{equation}
where
\begin{equation}
D_n = \frac{2}{L} \int_0^L f(x) \sin \left(\frac{n \pi x}{L} \right) dx
\end{equation}
We can compare these analytical results to the numerical ones obtained through finite differences, this is done in Figure ~\ref{fig:smoluchowskiCompare}

\begin{figure}
	\begin{subfigure}{0.49\textwidth}
	\includegraphics[width=\columnwidth]{temperatureAnalytic}
	\end{subfigure}
	\begin{subfigure}{0.49\textwidth}
	\includegraphics[width=\columnwidth]{temperatureAbsDiff}
	\end{subfigure}
	\caption{Finite differences on the heat equation: We simulate a situation where there are no sources for the heat equation and use equation \ref{eqn:analTemperature} to obtain the analytical result we then compare this to the result of finite differences with $\Delta t = 10^{-2}$ and $\Delta x = 0.02$. \label{fig:temperatureCompare}}
\end{figure}

\subsection{Convergence tests}
The idea of finite differences is that as the step size goes to zero, the numerical approximation will converge on the correct solution to the underlying equation being approximated. In the previous section we showed that finite differences approximated the solution very closely in some instances where the analytical result could be obtained. In general, we will not have an analytical solution to compare to but we would still like to be able to quantify the performance of our techniques.

Convergence tests involve decreasing the discretization size and checking whether the numerical solutions converge at all. In Figure \ref{fig:Convergence}, the Smoluchowski equation was simulated while keeping the temperature fixed, each time we halve $\Delta t$ and measure the normed difference between the new result and the previous one.

%Here we will decrease the step size and see whether or not the numerical scheme converges on a particular result

\begin{figure}[tb]
	\begin{subfigure}{0.49\textwidth}
		\includegraphics[width=\columnwidth]{probabilityConvergence}
	\end{subfigure}
	\begin{subfigure}{0.49\textwidth}
		\includegraphics[width=\columnwidth]{probabilityConvergenceRate}
	\end{subfigure}
\caption{The convergence of the probability distribution as $\Delta t$ is decreased, the spatial discretization $\Delta x$ is kept constant at 0.006. (a) the Smoluchoski equation is simulated forward for 1.0 seconds each line shows the result with a different value of the time step $\Delta t$. (b) the normed difference between each of the successive vectors is calculated and the result is plotted on  a log scale, the slope of this graph is called the convergence rate.}
\label{fig:Convergence}
\end{figure}

Likewise, we can do convergence tests for the coupled system, for brevity we have only included the results for the evolution of the temperature. All of these plots show that the approximation converges exponentially to the solution.

\begin{figure}[b]
	\begin{subfigure}{0.49\textwidth}
		\includegraphics[width=\columnwidth]{temperatureConvergence}
	\end{subfigure}
	\begin{subfigure}{0.49\textwidth}
		\includegraphics[width=\columnwidth]{temperatureNormDiff}
	\end{subfigure}
\caption{The convergence of the temperature as $\Delta t$ is decreased, the spatial discretization $\Delta x$ is kept constant at 0.006. (a) the coupled equations are simulated forward for 1.0 seconds each line shows the temperature with a different value of the time step $\Delta t$. (b) the normed difference between each of the successive vectors is calculated and the result is plotted on  a log scale.}
\label{fig:temperatureConvergence}
\end{figure}

%----------------------------------------------------------------------------------------
%	EXPLORATION
%----------------------------------------------------------------------------------------
\chapter{Exploration}

\setlength\epigraphwidth{.9\textwidth}
\epigraph{In that Empire, the Art of Cartography attained such Perfection that the map of a single Province occupied the entirety of a City, and the map of the Empire, the entirety of a Province. In time, those Unconscionable Maps no longer satisfied, and the Cartographers Guilds struck a Map of the Empire whose size was that of the Empire, and which coincided point for point with it. The following Generations, who were not so fond of the Study of Cartography as their Forebears had been, saw that that vast map was Useless, and not without some Pitilessness was it, that they delivered it up to the Inclemencies of Sun and Winters. In the Deserts of the West, still today, there are Tattered Ruins of that Map, inhabited by Animals and Beggars; in all the Land there is no other Relic of the Disciplines of Geography.}{\textbf{Jorge Luis Borges and Adolfo Bioy Cesares \\ -- ``On Exactitude in Science''}}

Now that we have set up the framework for our system, we are able to discuss particular phenomena and their relation to nature. Here we will discuss some prototypical problems in stochastic systems and their behavior when considered in the light of equations \ref{eqn:Smoluchowski} and \ref{eqn:TemperatureEvolution}. The point of this chapter is not to explain a particular experiment or the nature of a particular molecular motor in great detail. Nor is our aim to make a general model that may be easily adapted to any situation that the reader might fancy, but instead we aim to extract the physics that is at the essence of equations \ref{eqn:Smoluchowski} and \ref{eqn:TemperatureEvolution}. We will find that our system differs from those previously explored, however the previous behavior may always be recovered by taking the limit of either high specific heat capacity or fast thermal diffusivity as discussed in ~\autoref{dimensionless}.

\section{Bistable potentials} \label{Kramers}

\begin{figure}[tb]
\includegraphics[width=\columnwidth]{bistablePotential}
\caption{\textbf{Bistable potential.} In this plot we show the potential used to explore the Kramers rate and the reverse Landauer blowtorch, the potential has local minima at $a$ and $c$ and a maximum at $b$. If we begin with a probability distribution in the upper well, then the distribution will decay into the ground state of the upper well and then begin to decay into the lower well. The rate of flow from the upper well to the lower one will be denoted by $\kappa_+$ and the rate of flow from the lower well into the upper one will be denoted by $\kappa_-$. In order for a Brownian particle to go from $a$ to $c$, it will need to acquire thermal energy $E_B^+$, likewise to go from $c$ to $a$, it will need energy $E_B^-$.}
\label{fig:bistablePotential}
\end{figure}

A bistable potential has two stable minima and an intermediate unstable maximum as depicted schematically in Figure \ref{fig:bistablePotential}, these potentials occur in a wide range of applications including digital logic \cite{MyersCelebranoKrishnan2015,Landauer1961}, protein folding \cite{BryngelsonWolynes1989} and chemical reactions \cite{BernePecora1976}. The bistable potential well is one of the simplest ways to approach the Kramers' rate and many other important properties of a stochastic system \cite{MyersCelebranoKrishnan2015,Barcilon1996,SantamariaHolekGadomskiRubi2011}. In the context of Brownian motion, understanding the nature of bistable potentials can help one to understand more complicated potentials comprised of multiple deep wells.

\subsection{Kramer's rate} \label{kramersRate}
Consider the potential shown in Figure \ref{fig:bistablePotential}, if we begin in a state where we are certain that the particle is in the upper well, then as time passes, we should expect the probability distribution to move from point $a$ over the barrier at $b$ and into the well at point $c$. We will consider the regime where $E^+_B = V(x_b) - V(x_a) \gg k_B T$, where $T$ is constant, in this regime the rate at which the particles flow from $a$ to $c$ is given by the Eyring-Kramers' law \cite{Eyring1935, Kramers1940}.

We will now derive the Kramers' rate analytically for a constant temperature, an example of this derivation involving a slightly different Fokker-Planck equation is shown in Ref \cite{Gardiner2009}. Say that the particle begins at $a$, we are interested in knowing how long it will take to reach a point $x$. This function is denoted by $\tau(x)$ and is called the first passage time \cite{Gardiner2009} the first passage time is associated with a density $f(x, \, t)$. In particular, $f(x, \, t)$ denotes the probability that the particle has not passed the point $x$ after a time $t$, given that the particle started at $c$. We would like to obtain the mean first passage time $\langle \tau(x) \rangle \equiv \int_0^{\infty} t \, f(x, \, t) \, dt$, to do this we consider the dimensionless form of the Smoluchowski equation: 

\begin{equation}
\frac{\partial P(x, \, t)}{\partial t} = \frac{\partial}{\partial x} \left(P(x, \, t) \partial_x V + T \partial_x P(x, \, t) \right)
\end{equation}
we see that $f(x, \, t)$ obeys the same equation of motion as $P(x, \, t)$. Integrating both sides yields:

\begin{equation}
\int_0^{\infty} dt \partial_t f(x, \, t) = \int_0^{\infty} dt \frac{\partial}{\partial x} \left ( f(x, \, t) \partial_x V + T \partial_x f(x, \, t) \right)
\end{equation}

using the fact that $f(x, \, \infty) = 0$ and $f(x, \, 0) = 1$, we find that:

\begin{equation}
-1 = \frac{\partial}{\partial x} \left (-\langle \tau(x) \rangle \partial_x V + T \frac{ \partial \langle \tau(x) \rangle}{\partial x} \right)
\end{equation}

Since time dependence has been removed and we are only in one dimension, this is an ordinary differential equation for $\langle \tau(x) \rangle$. We notice that,

\begin{equation}
\frac{d}{dx} \left [\exp \left (-\frac{V(x)}{T} \right) \frac{d}{dx} \langle \tau(x) \rangle \right] = \frac{1}{T} \exp \left (\frac{-V(x)}{T} \right ) \left (- V'(x) \frac{d}{dx} \langle \tau(x) \rangle + \frac{d^2}{dx^2} \langle \tau(x) \rangle \right)
\end{equation}

meaning that

\begin{equation}
\frac{d}{dx} \left [\exp \left (-\frac{V(x)}{T} \right) \frac{d}{dx} \langle \tau(x) \rangle \right] = -\frac{1}{T} \exp \left(-\frac{V(x)}{T} \right)
\end{equation}
Integrating from $-\infty$ to $x$ we have:

\begin{equation}
\exp \left (-\frac{V(x)}{T} \right) \frac{d}{dx} \langle \tau(x) \rangle = -\frac{1}{T} \int_{-\infty}^x dz \exp \left(-\frac{V(z)}{T} \right)
\end{equation}

Finally, integrating from $a$ to $x$ and using the fact that $\langle \tau(a) \rangle = 0$, we get:

\begin{equation}
\langle \tau(x) \rangle = - \frac{1}{T} \int_a^{x} dy \exp \left(\frac{V(y)}{T} \right) \int_{-\infty}^y dz \exp \left(-\frac{V(z)}{T} \right)
\end{equation}

If the particle starts at $a$, then the potential in the first integrand can be approximated by the second order Taylor expansion around $b$ and the second integrand will be approximated by an expansion around $a$, more explicitly:

\begin{eqnarray}
V(y) &\approx& V(b) + \frac{V''(b)}{2}(x - b)^2 \\
V(z) &\approx& V(a) + \frac{V''(a)}{2}(x - a)^2 
\end{eqnarray}
So for the mean first passage time of particles going from $a$ to $c$ we find,

\begin{equation}
\langle \tau_{a \to c} \rangle = \frac{2 \pi}{\sqrt{-V''(b) V''(a)}} \exp{\left (\frac{V(b) - V(a)}{T} \right)}  \label{eqn:kramersRate}
\end{equation}


The Kramer's rate is given by one over the mean first passage time, we will denote the rate of particles moving from $a$ to $c$ as $\kappa_+$, which is given by:

\begin{equation}
\kappa_+ = \frac{\sqrt{|V''(b) V''(a) |}}{2 \pi} e^\frac{-E_B}{T}
\end{equation}


Likewise, there will be a current flowing from $c$ to $a$, we will denote this by $\kappa_-$, once we have calculated both of these rates, the population in the upper well will be given by the differential equation

\begin{equation}
\frac{d P_+}{dt} = -\kappa_+ P_+(t) + \kappa_- P_-(t)
\end{equation}

If we start with the population situated entirely in the upper well, then we get:

\begin{equation}
P_+(t) = \exp{(-(\kappa_+ - \kappa_- )t)}.
\end{equation}

\begin{figure}
	\center
	\includegraphics[width=\textwidth]{kramersLinear}
	\caption{\textbf{The probability of being to the right of the barrier.} Here we see that the system decays into the ground state of the upper well and then begins to decay exponentially into the lower well. At the end of this, the probability of finding the particle on the upper well is almost zero. We can fit an exponential to this decay in order to estimate the Kramers' rate. \label{fig:kramersLinear}}
\end{figure}

We can also achieve this result numerically by starting the system off in the upper well and simulating forward while calculating the probability that the particle is in the upper well at each step. We then fit an exponential to this data and the fitted rate will be our numerically estimated Kramers rate. To assist with measuring the Kramer's rate, we used Hermite interpolation to create a sixth order polynomial with the desired bistable shape. Specifically, we would allow $V(x)$ to be a general 6th order polynomial and then we would specify the value of the polynomial at the locations $a$, $b$ and $c$, we would also enforce that the first derivative vanished at these locations. This would yield 6 equations which we would solve to give us the coefficients of $V(x)$, an example of such a polynomial is shown in Figure \ref{fig:bistablePotential}. Using these interpolating polynomials, we could keep the locations of the wells fixed while controlling the barrier height $E_B^+$. Once we had a potential, we would place the probability distribution in the upper well and then measure the probability that the particle is in the upper well with time and use this information to obtain the Kramer's rate, as shown in Figure \ref{fig:kramersLinear}.

This technique was used to calculate the Kramer's rate for multiple barrier height and potentials and it was found to be in excellent agreement with equation \ref{eqn:kramersRate}. The same technique can be used to obtain the Kramer's rate when the temperature is not held fixed. In this case, the Kramer's rate is sensitive to the boundary conditions imposed on the temperature and on the precise values of $\alpha$ and $\beta$ as shown in Figure \ref{fig:kramersAlphaBeta}.

\begin{figure}
	\center
	\includegraphics[width=\textwidth]{kramersAlphaBeta}
	\caption{\textbf{Kramers' rate} Here we vary $\alpha$ and $\beta$ and perform the same measurements as in Figure \ref{fig:kramersLinear}, the results show that the Kramers' rate varies with $\alpha$ and $\beta$. \label{fig:kramersAlphaBeta}}
\end{figure}

We therefore conclude that the Kramer's rate is dependent on the heat capacity and thermal diffusivity of the system.

\subsection{The reverse Landauer blowtorch} \label{reverseLandauerBlowtorch}

\begin{figure}
	\begin{subfigure}{0.49\textwidth}
		\includegraphics[width=\textwidth]{reverseBlowtorchInit}
	\end{subfigure}
	\begin{subfigure}{0.49\textwidth}
		\includegraphics[width=\textwidth]{reverseBlowtorchFinal}
	\end{subfigure}
	\caption{\textbf{Reverse Landauer blowtorch effect.} (a) The system starts off in the upper well of a bistable potential with a uniform temperature, we have $\alpha = 8 \cdot 10^{-2}$ and $\beta = 1.5$. (b) As the system decays into the lower well, the bath loses thermal energy in the form of heat, thus we see Brownian cooling as a consequence of the forcing of the potential. In these simulations, Dirichlet boundary conditions were imposed, meaning that the temperature at the ends was held fixed. This means that there was heat flowing through the boundaries, this heat was calculated at each time step and used to update the energy at each time step. This is necessary because the numerics use the energy to normalize the temperature as mentioned in ~\autoref{numerics}. \label{fig:reverseBlowtorch}}
\end{figure}
As noted in ~\autoref{landauersBlowtorch}, the relative occupancy of wells depends on the spatial distribution of the temperature, meaning that the temperature can act as a pseudo-force. However, it has been noted that when the movement of the particle has an effect on the environment, the opposite can occur \cite{DasDasBarikEtAl2015}. Concretely, the the reverse Landauer blowtorch effect says that if there is a force applied to a Brownian particle, then the force will cause the particle to induce a temperature gradient in the environment. This is shown in Figure \ref{fig:reverseBlowtorch}, in this figure, the temperature is held fixed at the boundaries, which we interpret physically as meaning that the domain is embedded in a much larger system that is held at a fixed temperature. Another example of a stochastic process that has an effect on the temperature of the environment is the Soret-Dufour effect which has been observed experimentally \cite{Onsager1931,HortLinzLuecke1992,PiazzaGuarino2002}, in the Soret effect or the Dufour effect, a non-equilibrium concentration of molecules is capable of producing a temperature gradient. Ref \cite{Santamaria-HolekGadomskiRubi2011} showed that thermal gradients due to the Soret effect have to be taken into account to explain measurements performed on protein crystal growth. Despite the knowledge of these self-induced temperature gradients, to our knowledge, Figure \ref{fig:reverseBlowtorch} is the first concrete demonstration of Brownian cooling as described by a self consistent theoretical model.

%they have moved into the lower well they will cause the upper well to become colder in a process called Brownian cooling, a simulation of this process is shown in Figure \ref{fig:reverseBlowtorch}. In this paper, the authors treat the system stochastically and are not able to model the diffusion of temperature. This is the same as our model with $\beta$ in equation \ref{eqn:dimensionlessHeat} set to zero. 

\section{Titlted periodic potentials}
Tilted periodic potentials are very important in biology where they can be used to model molecular motors, in Ref \cite{MaLaiAckersonEtAl2015, MaLaiAckersonEtAl2015a} the authors synthesize a tilted periodic potential by placing Brownian particles on a crystalline surface. The surface is then titled at an angle $\theta$ relative to the normal defined by gravity, in their theoretical analysis, the authors assume a constant temperature and are therefore able to describe a non-equilibrium steady state. Interestingly, in section \ref{SteadyState}, we showed that our system does not have such a well defined steady state for tilted periodic potentials.

Despite not being able to easily calculate the steady state for titled periodic potentials, we can still explore the dynamical case as shown in Figure \ref{fig:tiltedPeriodic}. The tilted periodic potential combines the physics of both the Kramers' rate and the reverse Landauer blowtorch as described in ~\autoref{Kramers}. One can imagine the potential in Figure \ref{fig:tiltedPeriodic} as being a multi-stable potential as opposed to the bistable potential shown in Figure \ref{fig:bistablePotential}. In fact, noticing that the titled periodic potential is a multi-stable potential led the authors of Ref \cite{ChallisJack2014} to describe a periodic potential with a master equation that involved the Kramers' rate of each of the adjacent wells. As well as the return of the Kramers' rate, Figure \ref{fig:tiltedPeriodic} is a very good demonstration of the reverse Landauer blowtorch effect. One can see form looking at the figure that the Brownian particle is absorbing heat from the environment as it moves down the potential. We therefore speculate that the reverse Landauer blowtorch effect could have an effect on biological molecular motors.

\begin{figure}
	\begin{subfigure}{0.49\textwidth}
		\includegraphics[width=\textwidth]{tiltedPeriodicInit}
	\end{subfigure}
	\begin{subfigure}{0.49\textwidth}
		\includegraphics[width=\textwidth]{tiltedPeriodicFinal}
	\end{subfigure}
	\caption{\textbf{Tilted periodic potential.} Imposing Dirichlet boundary conditions maintaining that the temperature is equal to 0.8 at the boundaries (temperature is dimensionless), we have $\alpha = 1 \cdot 10^{-4}$ and $\beta = 1 \cdot 10^{-2}$. (a) The system begins with a uniform temperature with a Gaussian probability distribution located near the top of the potential. (b) after some time, the particle has moved down the potential, removing heat from the environment as it overcomes the periodic barriers provided by the potential. We notice that unlike Figure \ref{fig:Schematic}, the probability distribution does not decay into Gaussian distributions in the wells, instead the temperature gradients cause the probability distribution to take on a very complicated form. \label{fig:tiltedPeriodic}}
\end{figure}

In order to quantify the effect of $\alpha$ and $\beta$ on titled periodic wells we will turn to potentials that have deep wells, in these potentials the probability distribution will be very strongly localized at the bottom of the wells and will be hopping from well to well as explained in ~\autoref{kramersRate}. As the particle hops between wells, it will reduce the temperature in the well that it is currently situated in. When $\alpha$ is made large enough, then the Landauer blowtorch effect may cause the particle to hop back into the well that it came from. This is shown in Figure \ref{fig:tiltedPeriodicQuant} and Figure \ref{fig:alphaCritical}, here we see that there is a critical value of $\alpha$.
\begin{figure}
	\begin{subfigure}{0.49\textwidth}
		\includegraphics[width=\textwidth]{tiltedPeriodicQuantInit}
	\end{subfigure}
	\begin{subfigure}{0.49\textwidth}
		\includegraphics[width=\textwidth]{tiltedPeriodicQuantFinal}
	\end{subfigure}
	\begin{subfigure}{0.49\textwidth}
		\includegraphics[width=\textwidth]{tiltedPeriodicQuantWellTemps}
	\end{subfigure}
	\begin{subfigure}{0.49\textwidth}
		\includegraphics[width=\textwidth]{tiltedPeriodicQuantWellProbs}
	\end{subfigure}
	\caption{\textbf{Stalling in the tilted periodic potential} $\alpha = 8.3 \times 10^{-5}$ and $\beta = 1 \times 10^{-2}$ (a) \label{fig:tiltedPeriodicQuant}}
\end{figure}
Below this critical value, the self-induced temperature gradient will not be substantial so the particle will hop out of the higher well and drift down the potential. However, if $\alpha$ is above the critical value, then the self-induced temperature gradient caused by the movement of the particle will cause the particle to actually move back into the well that it came from, at the same time causing the temperature of that well to decrease. Eventually the particle reaches a steady state where the temperature gradient causing the particle to move against the potential is exactly canceled by the force of the potential itself. When one is at the critical value of $\alpha$, the probability distribution does not fully leak out of the initial well or hop back into it and instead will be constantly balancing between these two limits. This critical value is highly sensitive to $\alpha$, the probability may remain flat for quite some time only to either hop down to the lower well or hop back up into the higher well. In the context of molecular motors, the critical value of $\alpha$ would correspond to a stall.

\begin{figure}
	\raggedleft
	\includegraphics[width=0.9\textwidth]{probabilityInitialWell}
	\caption{\textbf{The population of the initial well for different values of $\alpha$} \label{fig:alphaCritical}}
\end{figure}

%Consider the case where $\alpha$ is above it's critical value, at this point the probability distribution will reach a steady state where the probability distribution is located at two wells. We speculate that this may be a reason why equations \ref{eqn:Smoluchowski} and \ref{eqn:TemperatureEvolution} do not admit a unique steady state for tilted potentials. To see this, note that the initial location of the probability distribution was chosen at will, had we chosen a different location, then the steady state distribution that the particle eventually attains would be different. At this stage we are not sure whether or not the steady state observed is a quasi-static one or a truly steady state. A quasi-steady state is a state where the probability has not completely stopped evolving, but it has slowed down a lot. Figure \ref{fig:meSimulateYouLongTime} shows the result of a simulation where $\alpha$ was greater than the critical value and we continued observing the evolution of the particle for a very long time. From this, we are unsure whether or not this is a steady state or a quasi-steady state.

%----------------------------------------------------------------------------------------
%	DISCUSSION
%----------------------------------------------------------------------------------------
\chapter{Discussion}

We have shown that self induced temperature gradients can play a crucial role in Brownian dynamics. Purely from a theoretical stand-point, such a model better explains the flow of energy and entropy through a system. In ~\autoref{thermodynamics} we explained how the temperature of the environment is a crucial player in the calculation of the entropy and energy of the system. The fact that our system has a well defined entropy is a strong advantage and allows one to talk about information in a rigorous way \cite{Landauer1961,MyersCelebranoKrishnan2015}.

%----------------------------------------------------------------------------------------
%	OUTLOOK
%----------------------------------------------------------------------------------------
\chapter{Outlook} \label{outlook}

This chapter will discuss possible future work and goals that are beyond the scope of this honors project, the main focus of this section is as follows:

\begin{itemize}
    \item{Generalization to 2d and 3d systems}
    \item{Inter-particle interactions}
    \item{Flow of heat via convection as well as an extension to compressible fluids and gases}
\end{itemize}

\section{Generalization to 2d and 3d systems}
Many systems of interest are multi-dimensional and cannot be approximated by a one-dimensional system {\color{red} citations}. In particular, finding the efficiency of a molecular motor involves calculating the flow of energy from one degree of freedom to another {\color{red} [citations]}. Furthermore, generalizing our system to more dimensions gives more directions in which heat can flow, therefore allowing us to conserve energy in a more general sense {\color{red} is this true? We can sort of do things without this by considering energy flows at the 1d boundaries.}

\section{Inter-particle interactions}
Inter-particle interactions play a very important role in Brownian dynaimcs {\color{red} [citations]}, an example of this is the myosin protein in cells which is used in the contraction of muscles {\color{red} citations}. Individually, myosin proteins are poor at carrying loads when compared to their kinesin couterparts {\color{red} [citations]}, however their collective actions are very powerful due to the rich nature of their inter-motor interactions.

The set of equations that we have been using doe not apply to more than one Brownian particle without some considerable changes. This is due to very complicated inter-particle interactions that can occur when considering thermal fluctautions of the kind that we are focusing on. To visualize this intuilively, imagine two particles diffusing in their reaction coordinates, both subject to the same potential. Particle one may induce local temperature gradients that will affect particle two if particle two moves to where particle one created those temperature gradients. These interactions involve second order statistics that are beyond the scope of this project.

\section{Fluid dynamics}
The Brownian particles that we have been modelling have been suspended in a fliud. The dynamics of fluids present a formidablle and exciting challenge to both physicicsts and mathematicians alike {\color{red} [citations]}, as well as this there has been effort in the literature to combine the Smoluchowski equation with the Navier-Stokes equation, thus linking the worlds of Brownian dynamics with fluid dynamics {\color{red} [citations]}. The two physical phenomena that we neglect which could be treated with techniques from fluid dynamics are compressible fluids and heat convection through fluid flow. Despite the fact that we neglected these phenomena, our system is still physical and is self consistent.

In section \ref{thermodynamics}, in order to gain insight on the entropy we had to assume that we were dealing with an incompressible fluid. Without this assumption, we would not be able to gauruntee that our equations of motion
increase the entropy. Therefore, in order to obey the laws of thermodynamics in a compressible fluid, we would have to include extra terms in our equations of motion.

The second important phenomena that can be described by fluid dynamics is fluid flow. This phenomena is responsible for heat transport via convection, currently we only consider heat flow via diffusion, this means that we are tacitly assuming that the fluid in which our Brownian particle is suspended, is stationary. Including convection into our model would involve modifying our heat equation so that it reads {\color{red} include the modified heat equation}.


%----------------------------------------------------------------------------------------
%	BIBLIOGRAPHY
%----------------------------------------------------------------------------------------

%\renewcommand{\refname}{\spacedlowsmallcaps{References}} % For modifying the bibliography heading

\clearpage{}   % Make sure that the references are on a new page.
\bibliography{BibData}{} % The file containing the bibliography

\bibliographystyle{unsrt}

%----------------------------------------------------------------------------------------
%	APPENDIX
%----------------------------------------------------------------------------------------

\makeatletter\@openrightfalse
\begin{appendices}

\chapter{Additional figures}
%\renewcommand{\textflush}{flushright} \renewcommand{\sourceflush}{flushright}
\setlength\epigraphwidth{.365\textwidth}
\epigraph{I am a fish.}{\textbf{Dr. Philip Brydon -- 2016}}

There are a lot of interesting phenomena that occur when one explores the system that we explored in detail. In this appendix, we would like to share some of the exciting results that we found during the project, these results do not directly contribute to the narrative of the thesis which is why they were not included in the main body of the text. Despite not being a part of the main body, we feel that the reader would be at a loss without being made aware of these phenomena.

Figure \ref{fig:exciting} shows an example of a complicated periodic potential, here we see that the evolution of the system is highly non-linear.
\begin{figure}
	\begin{subfigure}{0.49\textwidth}
		\includegraphics[width=\textwidth]{excitingInit}
	\end{subfigure}
	\begin{subfigure}{0.49\textwidth}
		\includegraphics[width=\textwidth]{excitingFinal}
	\end{subfigure}
	\caption{\textbf{Periodic boundary conditions with a complex periodic potential} $\alpha = 5 \times 10^{-3}$ and $\beta = 1 \times 10^{-2}$. (a) shows the initial configuration (b) shows the system after 1.5 dimensionless time units. \label{fig:exciting}}
\end{figure}

\subsubsection{Stochastic methods}
As mentioned in the introduction the underlying nature of Brownian motion is stochastic and can be understood through the Langevin equation. We solved the Langevin equation using the Euler-Maryama method for stochastic equations. An example of this is shown in Figure \ref{fig:stochastic}, here we see that the stochastic approach truly does have the same physics as equation \ref{eqn:Smoluchowski}.


\begin{figure}
	\center
	\includegraphics[width=0.9\textwidth]{Stochastic}
	\caption{\textbf{A stochastic simulation of the Langevin equation} A simulation of 1000 particles on a tilted periodic potential, all particles began at position 0 and were simulated forward for 2 seconds with a time step of 0.002 seconds. \label{fig:stochastic}}
\end{figure}


\end{appendices}
\@openrighttrue\makeatletter
%------------------------------------------------

%------------------------------------------------

%----------------------------------------------------------------------------------------

\end{document}
